\documentclass[11pt]{article}
\usepackage[margin=1in]{geometry}
\usepackage{amsmath, amssymb, amsthm, enumitem}

\title{\textbf{GRE Mathematics Subject Test}\\[4pt] Linear Algebra Practice Problems}
\author{}
\date{}

\newcounter{prob}
\newcommand{\problem}{\stepcounter{prob}\medskip\noindent\textbf{Problem \theprob.}\ }
\DeclareMathOperator{\tr}{tr}
\DeclareMathOperator{\rank}{rank}
\DeclareMathOperator{\nullity}{nullity}
\DeclareMathOperator{\Span}{span}
\DeclareMathOperator{\diag}{diag}

\begin{document}
\maketitle
\thispagestyle{empty}

%======================================================================
\section*{A. Vector Spaces, Subspaces \& Dimension}
%======================================================================

\problem
What is $\dim(P_4(\mathbb{R}))$, the space of real polynomials of degree at most 4?

\problem
Let $W=\{(x,y,z)\in\mathbb{R}^3 : x+2y-z=0\}$. What is $\dim W$?

\problem
Which of the following are subspaces of $\mathbb{R}^3$? Select \textbf{all} that apply.
\begin{enumerate}[label=(\alph*)]
\item $\{(x,y,z) : x+y+z=0\}$
\item $\{(x,y,z) : x+y+z=1\}$
\item $\{(x,y,z) : xz=0\}$
\item $\{(x,y,0) : x,y\in\mathbb{R}\}$
\end{enumerate}

\problem
True or False: The union of two subspaces of $V$ is always a subspace of $V$.

\problem
Let $W_1$ and $W_2$ be subspaces of $\mathbb{R}^5$ with $\dim W_1 = 3$ and $\dim W_2 = 4$. What is the minimum possible dimension of $W_1\cap W_2$?

%======================================================================
\section*{B. Linear Transformations \& Rank--Nullity}
%======================================================================

\problem
Let $T\colon\mathbb{R}^4\to\mathbb{R}^3$ be a linear transformation with $\rank(T)=2$. What is $\nullity(T)$?

\problem
Let $T\colon\mathbb{R}^3\to\mathbb{R}^3$ be defined by $T(x,y,z)=(x+y, y+z, 0)$. Find $\dim(\ker T)$.

\problem
Let $T\colon V\to W$ be a linear transformation with $\dim V=6$ and $\dim W=4$. What is the maximum possible rank of $T$?

\problem
True or False: If $T\colon\mathbb{R}^5\to\mathbb{R}^3$ is a linear transformation, then $T$ cannot be injective.

\problem
Let $A$ be a $3\times 5$ matrix with $\rank(A)=3$. How many free variables does the system $Ax=0$ have?

%======================================================================
\section*{C. Matrices \& Determinants}
%======================================================================

\problem
Compute $\det\begin{pmatrix}2&1&0\\0&3&1\\0&0&4\end{pmatrix}$.

\problem
Let $A$ be a $4\times 4$ matrix with $\det A = 3$. What is $\det(2A)$?

\problem
If $A$ is a $3\times 3$ matrix with $\det A=5$, what is $\det(A^{-1}\cdot\text{adj}(A))$?

\problem
Let $A$ be a $3\times 3$ matrix with eigenvalues $1, 2, 3$. What is $\det(A^2+I)$?

\problem
True or False: If $AB=0$ for nonzero matrices $A$ and $B$, then $\det A = 0$.

%======================================================================
\section*{D. Eigenvalues \& Eigenvectors}
%======================================================================

\problem
Find the eigenvalues of $A=\begin{pmatrix}4&1\\0&3\end{pmatrix}$.

\problem
Let $A$ be a $3\times 3$ matrix with $\tr(A)=6$ and $\det(A)=6$. If the eigenvalues of $A$ are $1, 2, \lambda$, find $\lambda$.

\problem
Let $A$ be a $3\times 3$ matrix with eigenvalues $2, -1, 3$. What are the eigenvalues of $A^2-A$?

\problem
How many of the following matrices are diagonalizable over $\mathbb{R}$?
\begin{enumerate}[label=(\roman*)]
\item $\begin{pmatrix}1&0\\0&1\end{pmatrix}$
\item $\begin{pmatrix}1&1\\0&1\end{pmatrix}$
\item $\begin{pmatrix}0&-1\\1&0\end{pmatrix}$
\item $\begin{pmatrix}2&0\\0&3\end{pmatrix}$
\end{enumerate}

\problem
Let $A$ be a $4\times 4$ matrix with characteristic polynomial $(\lambda-1)^2(\lambda-3)^2$. If $\rank(A-I)=3$, what is the geometric multiplicity of $\lambda=1$?

%======================================================================
\section*{E. Special Matrices \& Properties}
%======================================================================

\problem
Let $A$ be a real symmetric $3\times 3$ matrix. Which of the following must be true? Select \textbf{all} that apply.
\begin{enumerate}[label=(\alph*)]
\item $A$ is diagonalizable.
\item All eigenvalues of $A$ are real.
\item Eigenvectors for distinct eigenvalues are orthogonal.
\item $A$ is invertible.
\end{enumerate}

\problem
Let $A$ be a nilpotent $3\times 3$ matrix. What is $\det(I+A)$?

\problem
If $A$ is idempotent (i.e., $A^2=A$) and $\rank(A)=3$ for a $5\times 5$ matrix $A$, what is $\tr(A)$?

\problem
Let $Q$ be a $3\times 3$ orthogonal matrix. Which of the following are possible values for $\det Q$?
\begin{enumerate}[label=(\alph*)]
\item $1$
\item $-1$
\item $0$
\item $2$
\end{enumerate}

%======================================================================
\section*{F. Inner Products, Orthogonality \& Applications}
%======================================================================

\problem
Let $W=\Span\{(1,1,0),(0,1,1)\}$ in $\mathbb{R}^3$. What is $\dim(W^\perp)$?

\problem
Apply one step of Gram--Schmidt: given $v_1=(1,1,0)$ and $v_2=(1,0,1)$, compute the orthogonal vector $u_2 = v_2 - \dfrac{\langle v_2,v_1\rangle}{\langle v_1,v_1\rangle}v_1$.

\problem
Let $A$ be a real $n\times n$ matrix. True or False: $\rank(A)=\rank(A^T A)$.

\problem
A real symmetric $3\times 3$ matrix $A$ has eigenvalues $2, 5, 5$. What is $\|A\|^2_F = \tr(A^TA)$?

\problem
Let $A=\begin{pmatrix}1&2\\2&4\end{pmatrix}$. Find the eigenvalues and determine whether $A$ is positive semi-definite.

\problem
Let $T\colon\mathbb{R}^3\to\mathbb{R}^3$ be the orthogonal projection onto $W=\Span\{(1,0,0),(0,1,0)\}$. What is the matrix of $T$ with respect to the standard basis?


%======================================================================
\newpage
\section*{Answer Key}
%======================================================================

\begin{enumerate}

\item $\dim(P_4(\mathbb{R}))=5$. A basis is $\{1, x, x^2, x^3, x^4\}$, which has 5 elements.

\item $\dim W=2$. $W$ is defined by one linear equation in $\mathbb{R}^3$, so $\dim W = 3-1=2$.

\item (a) and (d). (a): $x+y+z=0$ is a homogeneous linear equation, so the solution set is a subspace. (b): $0+0+0=0\neq 1$, so $\vec{0}\notin$ this set---not a subspace. (c): Not closed under addition: $(1,0,0)$ and $(0,0,1)$ satisfy $xz=0$, but $(1,0,1)$ does not. (d): This is the $xy$-plane $=\{(x,y,0)\}$, a subspace.

\item False. For example, in $\mathbb{R}^2$, let $W_1=\Span\{(1,0)\}$ (the $x$-axis) and $W_2=\Span\{(0,1)\}$ (the $y$-axis). Then $(1,0)\in W_1\cup W_2$ and $(0,1)\in W_1\cup W_2$, but $(1,1)=(1,0)+(0,1)\notin W_1\cup W_2$. So $W_1\cup W_2$ is not closed under addition.

\item $\dim(W_1\cap W_2)\ge 3+4-5=2$ (by the dimension formula $\dim(W_1+W_2)=\dim W_1+\dim W_2-\dim(W_1\cap W_2)$ and $\dim(W_1+W_2)\le 5$). The minimum is $\mathbf{2}$.

\item $\nullity(T)=\dim(\mathbb{R}^4)-\rank(T)=4-2=\mathbf{2}$ (Rank--Nullity Theorem).

\item $T(x,y,z)=(x+y,y+z,0)$. Solve $\ker T$: $x+y=0$ and $y+z=0$. From these: $x=-y$, $z=-y$. So $\ker T=\Span\{(-1,1,-1)\}$ and $\dim(\ker T)=\mathbf{1}$.

\item $\rank(T)\le\min(\dim V,\dim W)=\min(6,4)=\mathbf{4}$.

\item True. If $T$ were injective, $\nullity(T)=0$, so $\rank(T)=\dim(\mathbb{R}^5)=5$. But $\rank(T)\le\dim(\mathbb{R}^3)=3$, a contradiction.

\item Free variables $= n-\rank(A)=5-3=\mathbf{2}$.

\item $\det = 2\cdot 3\cdot 4=\mathbf{24}$ (product of diagonal entries for a triangular matrix).

\item $\det(2A)=2^4\cdot\det A=16\cdot 3=\mathbf{48}$ (since $\det(cA)=c^n\det A$ for an $n\times n$ matrix).

\item $\text{adj}(A)=\det(A)\cdot A^{-1}$, so $A^{-1}\cdot\text{adj}(A)=\det(A)\cdot A^{-2}$. For a $3\times 3$ matrix, $\det(cB)=c^3\det B$, so $\det(\det(A)\cdot A^{-2})=(\det A)^3\cdot(\det A)^{-2}=5^3\cdot 5^{-2}=\mathbf{5}$.

\item Eigenvalues of $A$ are $1,2,3$. Eigenvalues of $A^2+I$ are $1^2+1=2$, $2^2+1=5$, $3^2+1=10$. So $\det(A^2+I)=2\cdot 5\cdot 10=\mathbf{100}$.

\item True. If $AB=0$ and $B\neq 0$, then $B$ has a nonzero column $b$, and $Ab=0$ with $b\neq 0$. So $A$ has a nontrivial kernel, meaning $A$ is not invertible, so $\det A=0$.

\item Eigenvalues of an upper triangular matrix are the diagonal entries: $\lambda=\mathbf{4}$ and $\lambda=\mathbf{3}$.

\item $\tr(A)=1+2+\lambda=6\Rightarrow\lambda=3$. Check: $\det(A)=1\cdot 2\cdot 3=6$. \checkmark\ So $\lambda=\mathbf{3}$.

\item If $\lambda$ is an eigenvalue of $A$, then $\lambda^2-\lambda$ is an eigenvalue of $A^2-A$. So: $2^2-2=2$, $(-1)^2-(-1)=2$, $3^2-3=6$. The eigenvalues are $\mathbf{2, 2, 6}$.

\item \textbf{2} matrices are diagonalizable over $\mathbb{R}$. (i) $I_2$: already diagonal. \checkmark\ (ii) $\begin{pmatrix}1&1\\0&1\end{pmatrix}$: eigenvalue $\lambda=1$ with algebraic multiplicity 2, but $\rank(A-I)=1$, so geometric multiplicity is 1. NOT diagonalizable. (iii) $\begin{pmatrix}0&-1\\1&0\end{pmatrix}$: characteristic polynomial $\lambda^2+1=0$, no real roots. NOT diagonalizable over $\mathbb{R}$. (iv) Already diagonal. \checkmark.

\item Geometric multiplicity of $\lambda=1$ equals $\nullity(A-I)=4-\rank(A-I)=4-3=\mathbf{1}$. (Note: algebraic multiplicity is 2, and $1\le\text{geo}\le\text{alg}$ is satisfied.)

\item (a), (b), and (c). Real symmetric matrices have: all real eigenvalues (b), are orthogonally diagonalizable (a), and eigenvectors for distinct eigenvalues are orthogonal (c). (d) is false: the zero matrix is symmetric but not invertible.

\item $\det(I+A)=\mathbf{1}$. A nilpotent matrix has all eigenvalues equal to 0. Thus eigenvalues of $I+A$ are all equal to $1+0=1$. So $\det(I+A)=1^3=1$.

\item $\tr(A)=\mathbf{3}$. If $A$ is idempotent, eigenvalues are 0 or 1. $\tr(A)=$ sum of eigenvalues $=\rank(A)$ for idempotent matrices. So $\tr(A)=3$.

\item (a) and (b) only. $Q^TQ=I$ implies $(\det Q)^2=1$, so $\det Q=\pm 1$. Thus 0 and 2 are impossible.

\item $\dim(W^\perp)=3-\dim W$. The two vectors $(1,1,0)$ and $(0,1,1)$ are linearly independent, so $\dim W=2$. Thus $\dim(W^\perp)=3-2=\mathbf{1}$.

\item $\langle v_2,v_1\rangle = 1\cdot 1+0\cdot 1+1\cdot 0=1$. $\langle v_1,v_1\rangle=1+1+0=2$. So $u_2=v_2-\frac{1}{2}v_1=(1,0,1)-\frac{1}{2}(1,1,0)=\left(\frac{1}{2},-\frac{1}{2},1\right)$.

\item True. This is a standard result. The null spaces of $A$ and $A^TA$ are equal: $Ax=0\Rightarrow A^TAx=0$, and $A^TAx=0\Rightarrow x^TA^TAx=0\Rightarrow\|Ax\|^2=0\Rightarrow Ax=0$. Equal null spaces means equal nullities, and by rank--nullity (both have $n$ columns), equal ranks.

\item $\|A\|^2_F=\tr(A^TA)$. For a symmetric matrix, $A^TA=A^2$, so $\tr(A^2)=\sum\lambda_i^2=2^2+5^2+5^2=4+25+25=\mathbf{54}$.

\item Characteristic polynomial: $\det(A-\lambda I)=(1-\lambda)(4-\lambda)-4=\lambda^2-5\lambda=\lambda(\lambda-5)$. Eigenvalues: $\lambda=0$ and $\lambda=5$. Since $A$ is symmetric and all eigenvalues are $\ge 0$, $A$ is \textbf{positive semi-definite} (but not positive definite, since $0$ is an eigenvalue).

\item $T$ projects onto the $xy$-plane: $T(x,y,z)=(x,y,0)$. The matrix is $\begin{pmatrix}1&0&0\\0&1&0\\0&0&0\end{pmatrix}$.

\end{enumerate}

\end{document}
